% ===================== Phase 2: Implementation and evaluation =====================
% Every \textbf{[...]} is a placeholder for a number from YOUR runs (README, "Phase 2
% experiment"). Sandbox numbers from development are not valid report data.
% Figures and the latency table are produced by bench.plot; copy them into figures/.
\section{Phase 2: Implementation and Latency Evaluation}\label{sec:impl}

\subsection{Implementation}
Each component of Fig.~\ref{fig:arch} runs in its own container, defined in one
Compose file (versions and settings in Table~\ref{tab:setup}); the server is
limited to one CPU so that results do not depend on the host's core count. Each
text is tokenized once at startup with the word rule of FR1. A cache miss scans
the text's word sequence; we deliberately keep no precomputed word index, which
would leave the cache nothing to save. Requests are fully validated before Redis
is touched, so malformed requests never count as hot keywords; if Redis is down,
the server answers without it (fail-open). Every request opens a new connection,
because the Phase~3 load balancer forwards byte streams and can only choose a
server when a connection opens. Tracing the protocol messages shows the price:
besides the TCP handshake, a count takes four RPyC round trips (\textsc{getroot},
\textsc{inspect}, \textsc{getattr}, \textsc{call}) before the reply arrives,
where a reused connection needs two.

\subsection{Experiments}
Latency is measured at the client, from opening the connection (part of sending
the request) to receiving the reply. A seeded trace of 10{,}000 requests
(Table~\ref{tab:setup}) is replayed \emph{open-loop}: requests leave at their
scheduled times whether or not earlier replies have arrived, so the generator
cannot slow down with the server and hide queueing. Every run replays the same
trace from an empty cache, with inter-arrival gaps scaled to the rate, so only the
load differs between rates. Replies are checked against independently computed
counts, and a run is discarded if the generator lags its schedule (99th-percentile
send lag above 10\,ms). Rates were chosen after a calibration sweep, up to just
below the saturation point of one server (\textbf{[Rsat]}\,req/s).

\begin{table}[t]
\caption{Phase 2 experiment setup}
\label{tab:setup}
\centering
\footnotesize
\begin{tabular}{@{}p{0.27\columnwidth}p{0.67\columnwidth}@{}}
\hline
Host & \textbf{[CPU model, cores, RAM, OS]}; Docker \textbf{[version]}, \textbf{[n]} CPUs and \textbf{[m]}\,GB available to Docker \\
Software & Python 3.12, RPyC 6.0.2, redis-py 8.1.0; Redis 7.4, in-memory only, 256\,MB, LRU eviction \\
Containers & server: 1 CPU; Redis and client: no limit \\
Corpus & 8 Project Gutenberg books, licence text removed; \textbf{[N]} words \\
Trace & 10{,}000 requests, seed 42 \\
Texts & uniform over the 8 books \\
Keywords & pool of 1{,}000 words (corpus frequency ranks 100--5{,}099), Zipf popularity ($s=1.0$); best possible hit ratio \textbf{[X]}\% \\
Arrivals & Poisson, open loop; same trace at every rate \\
Rates & \textbf{[R1, R2, R3, R4, R5]} requests/s \\
Runs & 3 per rate, each from an empty cache \\
Metrics & mean and 99th-percentile latency (linear interpolation between ranks) \\
\hline
\end{tabular}
\end{table}

\subsection{Results}
Figs.~\ref{fig:mean} and~\ref{fig:p99} show the mean and the 99th-percentile
latency per rate; Table~\ref{tab:latency} lists the values. All
\textbf{[150{,}000]} requests were answered correctly and without errors, and
the hit ratio was \textbf{[X]}\% at every rate.
% TODO after the runs: replace every bracketed value and keep only the statements
% your figures support. The reasoning below is the expected explanation; check it.
At the lowest rate the mean latency is \textbf{[x]}\,ms. Most of it is a fixed
per-request cost, connection setup plus four RPyC round trips, which no cache hit
removes; a miss adds the scan of the requested text, longest for the largest book.
Up to \textbf{[R]}\,req/s the mean stays almost flat, while the 99th percentile
grows from \textbf{[a]} to \textbf{[b]}\,ms. The tail rises first because bursts
of arrivals queue behind slow requests: the server's CPU work is serialized, since
it has one CPU and Python's global interpreter lock runs one request thread at a
time. For such a single-server queue the mean waiting time grows with
$\lambda E[S^2]/(1-\rho)$ (Pollaczek--Khinchine~\cite{kleinrock1975queueing}),
where $\lambda$ is the arrival rate, $S$ the service time and $\rho$ the
utilization. Mixing fast hits with slow misses makes $E[S^2]$ large, so waiting
time, and the tail most of all~\cite{dean2013tail}, grows well before the server
is fully utilized. At \textbf{[R5]}\,req/s the 99th percentile reaches
\textbf{[c]}\,ms and violates NFR1: one server does not scale, which motivates
replicating it behind a load balancer (Phase~3).

\begin{figure}[t]
\centering
\IfFileExists{figures/latency_mean.pdf}{%
  \includegraphics[width=\columnwidth]{figures/latency_mean}}{%
  \fbox{\parbox[c][1.9in][c]{0.95\columnwidth}{\centering Run \texttt{bench.plot},
  then copy \texttt{latency\_mean.pdf} into \texttt{figures/}.}}}
\caption{Mean latency per request rate, one server (mean of 3 runs; whiskers: min--max).}
\label{fig:mean}
\end{figure}

\begin{figure}[t]
\centering
\IfFileExists{figures/latency_p99.pdf}{%
  \includegraphics[width=\columnwidth]{figures/latency_p99}}{%
  \fbox{\parbox[c][1.9in][c]{0.95\columnwidth}{\centering Run \texttt{bench.plot},
  then copy \texttt{latency\_p99.pdf} into \texttt{figures/}.}}}
\caption{99th-percentile (tail) latency per request rate, one server (mean of 3 runs; whiskers: min--max).}
\label{fig:p99}
\end{figure}

\begin{table}[t]
\caption{Latency per request rate, one server (ms; mean of 3 runs)}
\label{tab:latency}
\centering
\footnotesize
\IfFileExists{figures/latency_table.tex}{\input{figures/latency_table.tex}}{%
  \fbox{Run \texttt{bench.plot}, then copy \texttt{latency\_table.tex} into \texttt{figures/}.}}
\end{table}
